\begin{abstract}
\noindent
We develop a multi-layer stochastic structural model for valuing publicly traded companies whose primary asset is Bitcoin and that employ a leverage flywheel to grow per-share BTC exposure. The canonical example is Strategy (formerly MicroStrategy, MSTR), which as of late 2025 holds over 649{,}000 BTC financed through a complex capital stack of common equity, convertible bonds, and multiple series of perpetual preferred stock. No formal structural framework exists for these ``Bitcoin treasury vehicles,'' leaving investors without rigorous tools to assess whether equity premia reflect achievable BTC accretion or unsustainable leverage.

Our model couples four stochastic processes under the physical measure: (i) BTC price dynamics with jump-diffusion, (ii) an Ornstein--Uhlenbeck process for the equity-over-NAV log-premium capturing market sentiment and mean-reversion, (iii) a compound Poisson process for BTC holding accumulation driven by the leverage flywheel, and (iv) share-count dynamics from equity issuance and convertible conversion. We derive five novel diagnostic indicators---the Implied BTC Growth Rate (IBGR), Implied Leverage Elasticity (ILE), Total Equity Elasticity (TEE), Premium Mean-Reversion Index (PMRI), and Implied Funding Requirement Distribution (IFRD)---that decompose valuation into structural balance-sheet leverage, sentiment premium, and accretion efficiency.

Calibrating to Strategy's publicly available data through December 2025, we find: the model implies an 18.3\% annualized BTC accretion rate (18.4\% per-share), a structural leverage elasticity of 1.17, and a three-year survival probability exceeding 99.4\%. The current premium sits 1.5 standard deviations below its long-run mean, suggesting mean-reversion potential. We extend the framework conceptually to accommodate perpetual preferred stock layers (STRK, STRF) via a priority waterfall, laying groundwork for a complete multi-tranche capital structure analysis.
\end{abstract}
